% GitHub Repo and Documentation: https://github.com/celiobjunior/resume-template
% Copyright © 2025 Celio B Junior. All rights reserved.
% 
% Licensed under the Apache License, Version 2.0 (the "License");
% you may not use this file except in compliance with the License.
% You may obtain a copy of the License at
%
%     http://www.apache.org/licenses/LICENSE-2.0
%
% This template follows best practices from README.md

% Start of the LaTeX document: defines document type and format
% a4paper = A4 size, 10pt = base font size
\documentclass[a4paper,10pt]{article}

% --- PACKAGES: BASE / LANGUAGE ---
\usepackage[utf8]{inputenc}
\usepackage[T1]{fontenc}
\usepackage[english]{babel}

% --- PACKAGES: LAYOUT / TYPOGRAPHY ---
\usepackage{geometry}
\usepackage{parskip}
\usepackage{microtype}
\usepackage{enumitem}
\usepackage{titlesec}
\usepackage{array}
\usepackage{tabularx}

% --- PACKAGES: LINKS / BOOKMARKS ---
\usepackage{hyperref}
\usepackage{bookmark}
\usepackage{xurl}

% --- SETTINGS: MARGINS / SPACING ---
\geometry{top=1.5cm, bottom=1.5cm, left=1.5cm, right=1.5cm} % Page margins
\setcounter{secnumdepth}{0}                                 % Removes section numbering
\setlist[itemize]{
    leftmargin=0.75em, % List indentation (closer to page margin)
    itemsep=0.2em,     % Space between items
    topsep=0.25em,     % Space before/after list
    parsep=0em,        % Space between paragraphs in same item
    partopsep=0em      % Extra spacing at list start
}

% Removes page numbers and headers for a clean resume layout
\pagestyle{empty}

% --- SETTINGS: METADATA / LINKS ---
\hypersetup{
    pdftitle={CV Your Name},   % PDF title
    pdfauthor={Your Name},     % PDF author
    colorlinks=true,           % Uses colored links (no boxes)
    linkcolor=black,           % Internal links color
    urlcolor=black,            % URLs color
    citecolor=black,           % Citations color
    bookmarksdepth=2           % Sidebar: section and subsection
}

% --- SETTINGS: TITLES ---
\titleformat{\section}
{\Large\bfseries}
{}
{0em}
{}
[\titlerule\vspace{0.5ex}]
\titleformat{\subsection}
{\normalsize\bfseries}
{}
{0em}
{}
\titlespacing*{\subsection}{0pt}{0.35em}{0.15em}

% Macro for resume entries with PDF bookmark
\newcounter{cventry}
\newcommand{\cventry}[4]{%
    \refstepcounter{cventry}%
    \phantomsection%
    \pdfbookmark[2]{#1}{cventry-\thecventry}%
    \noindent\begin{tabularx}{\textwidth}{@{}>{\raggedright\arraybackslash}X >{\raggedleft\arraybackslash}X@{}}
    \textbf{#1} & #2 \\
    \textit{#3} & \textit{#4} \\
    \end{tabularx}
}

% --- START OF DOCUMENT ---
\begin{document}

% --- HEADER ---
% Replace with your personal information
\begin{center}
    {\LARGE \textbf{Full Name}} 
    \\ [0.1cm]
    City, State
    {\textbullet}
    \href{mailto:your.email@example.com}{your.email@example.com} 
    {\textbullet}
    \href{https://linkedin.com/in/username}{linkedin.com/in/username} 
    {\textbullet}
    \href{https://github.com/username}{github.com/username}
\end{center}

% --- SECTIONS ---

\section{Leadership Activities}
    % Include volunteer work, academic clubs, or leadership positions
    \cventry{Company/Organization Name}{Location (e.g., Remote or State, Country)}{Your Role/Position}{Start Date - End Date (e.g., Jan 2024 - Present)}
        \begin{itemize}
            % Use action verbs from README - quantify your impact whenever possible
            \item Describe a meaningful leadership achievement with specific metrics.
            \item Highlight how you managed teams, led initiatives, or drove outcomes.
            \item Show impact on the organization or community you served.
        \end{itemize}

% ------

\section{Experience}
    % List your professional or academic experiences, most recent first
    \cventry{Company/Organization Name}{Location (e.g., Remote or State, Country)}{Role (e.g., Software Engineer, Data Analyst)}{Start Date - End Date}
        \begin{itemize}
            % Focus on achievements, not only technical details - use the STAR method (Situation, Task, Action, Result)
            \item Achieved [specific result] through [action taken], resulting in [quantified impact].
            \item Developed/Built/Created [what you did] using [technologies/methods], resulting in [benefit].
            \item Collaborated with [team size/type] to [accomplish what], improving [metric] by [percentage/amount].
        \end{itemize}

    \cventry{Company/Organization Name}{Location (e.g., Remote or State, Country)}{Role (e.g., Junior Developer, Research Assistant)}{Start Date - End Date}
        \begin{itemize}
            \item Use active verbs from the action verbs table in README.md
            \item Quantify your achievements whenever possible (numbers, percentages, deadlines)
            \item Show progression and increasing responsibility across roles
        \end{itemize}

    \cventry{University/Research Institution}{Location (e.g., Remote or State, Country)}{Role (e.g., Undergraduate Researcher, Teaching Assistant)}{Start Date - End Date}
        \begin{itemize}
            \item Include research, academic work, or projects that demonstrate technical skills
            \item Highlight publications, presentations, or meaningful findings
        \end{itemize}

% ------

\section{Skills}
    % Keep this section concise and use as many keywords as possible
    % that are relevant to the role you want.
    \begin{itemize}
        \item \textbf{Technical:} List your programming languages, frameworks, tools, and methodologies
        \item \textbf{Languages:} Portuguese (Native), English (Proficiency Level), Third Language (Proficiency Level)
    \end{itemize}

% ------

% !!!!!!!!!!!!!!!!!!!!!!
% Move this section to the beginning, after the header and before
% your experiences, if you are looking for your first
% job or an internship. Otherwise, leave it as is.
% !!!!!!!!!!!!!!!!!!!!!!

\section{Education}
    % Most recent education first
    \cventry{Your University Name}{Location (e.g., Remote or State, Country)}{Your Degree Title (e.g., B.S. in Computer Science)}{Start Date - End Date (e.g., Jan 2020 - Dec 2024)}

\end{document}